% ===========================================================================
\chapter{Especificação dos Casos de Uso de Avaliação Docente e Indicadores}
\label{cap_avaliacao_indicadores}
% ===========================================================================

O fechamento do ciclo pedagógico e científico do SIGEA-GTT ocorre no subsistema de Avaliação Docente e Gestão de Indicadores Epidemiológicos. Neste estágio, o professor titular confronta a submissão discente com o gabarito clínico de referência, avalia a coerência da análise causal, atribui nota formativa com parecer pedagógico e gera os indicadores hospitalares padronizados pela metodologia IHI.

A Figura \ref{fig_uc_modulo_avaliacao} exibe o diagrama de casos de uso correspondente em notação UML e renderização nativa TikZ:

\begin{figure}[!ht]
\centering
\caption{Diagrama de Casos de Uso do Subsistema de Avaliação e Indicadores}
\label{fig_uc_modulo_avaliacao}
\resizebox{0.90\textwidth}{!}{%
\begin{tikzpicture}[
    actor/.style={
        draw=clinical-900,
        fill=clinical-100!40,
        thick,
        rounded corners=2pt,
        align=center,
        font=\bfseries\scriptsize,
        minimum width=2.4cm,
        minimum height=1.2cm
    },
    usecase/.style={
        ellipse,
        draw=clinical-700,
        fill=white,
        thick,
        align=center,
        font=\scriptsize,
        minimum width=3.4cm,
        minimum height=1.1cm,
        inner sep=2pt
    },
    link/.style={
        draw=clinical-900!80,
        thick
    },
    include_link/.style={
        draw=clinical-500,
        thick,
        dashed,
        -Latex
    }
]
    % Fronteira
    \draw[rounded corners=8pt, fill=clinical-100!15, draw=clinical-700, line width=1.2pt] 
        (0.5, -5.5) rectangle (12.5, 3.2);
    \node[anchor=north west, font=\bfseries\small\color{clinical-700}] at (0.8, 2.9) 
        {Subsistema: Avaliação Docente e Indicadores GTT};

    % Atores
    \node[actor] (prof) at (-2.2, 0.5) {Docente Titular\\\texttt{\scriptsize (PROFESSOR)}};
    \node[actor] (admin) at (-2.2, -3.5) {Administrador\\\texttt{\scriptsize (ADMIN)}};

    % Casos de Uso
    \node[usecase] (uc19) at (3.5, 1.8) {UC19: Avaliar Submissão\\e Emitir Feedback};
    \node[usecase] (uc20) at (8.5, 1.8) {UC20: Comparar com\\Gabarito Docente};
    \node[usecase] (uc21) at (3.5, -1.0) {UC21: Consolidar Indicadores\\Epidemiológicos GTT};
    \node[usecase] (uc22) at (8.5, -1.0) {UC22: Exportar Relatórios\\e Métricas Clínicas};

    % Conexões
    \draw[link] (prof) -- (uc19);
    \draw[link] (prof) -- (uc20);
    \draw[link] (prof) -- (uc21);
    \draw[link] (prof) -- (uc22);

    \draw[link] (admin) -- (uc21);
    \draw[link] (admin) -- (uc22);

    % Relação Include
    \draw[include_link] (uc19) -- node[above, font=\tiny\bfseries] {<<include>>} (uc20);

\end{tikzpicture}
}
\fonte{Elaborado pelo autor (2026).}
\end{figure}

% ===========================================================================
\section{Especificação Detalhada dos Casos de Uso}
\label{sec_especificacao_uc_avaliacao}
% ===========================================================================

% ---------------------------------------------------------------------------
% UC19
% ---------------------------------------------------------------------------
\subsection{UC19: Avaliar Submissão e Emitir Feedback Docente}
\label{uc19_especificacao}

\begin{table}[!ht]
\centering
\caption{Especificação Detalhada do Caso de Uso UC19}
\label{tab_uc19}
\footnotesize
\begin{tabularx}{\textwidth}{p{3.5cm} X}
\toprule
\textbf{Propriedade} & \textbf{Especificação Técnica e Comportamental} \\
\midrule
\textbf{Identificador e Nome} & \texttt{UC19: Avaliar Submissão e Emitir Feedback Docente} \\
\textbf{Atores} & \texttt{PROFESSOR} (Ator Primário). \\
\textbf{Nível de Meta} & Nível Usuário (\textit{Sea Level}). \\
\textbf{Pré-condições} & Existência de submissão de auditoria realizada pelo aluno na turma sob titularidade do professor. \\
\textbf{Gatilho} & O professor abre a lista de entregas da atividade e seleciona uma submissão pendente de nota. \\
\textbf{Fluxo Principal} & 
1. O sistema exibe o painel de correção contendo os gatilhos identificados pelo discente e os quatro artefatos da qualidade.\\
& 2. O sistema invoca automaticamente \texttt{UC20: Comparar Gatilhos com Gabarito Docente}, destacando em cores os acertos, omissões e falsos positivos do aluno.\\
& 3. O professor analisa a pertinência das justificativas clínicas e a qualidade do plano 5W2H/PDCA.\\
& 4. O professor atribui a nota numérica no campo específico.\\
& 5. O sistema valida se a nota informada está no intervalo válido de $0{,}00$ a $10{,}00$.\\
& 6. O professor digita o parecer pedagógico e formativo no campo \texttt{professor\_feedback}.\\
& 7. O professor clica em "Salvar Avaliação".\\
& 8. O sistema atualiza os atributos \texttt{grade}, \texttt{professor\_feedback} e \texttt{graded\_at} na tabela \texttt{activity\_submissions}.\\
& 9. A nota e o parecer tornam-se visíveis no painel individual do estudante avaliado. \\
\textbf{Fluxos de Exceção} & 
\textbf{[5e] Nota Fora da Faixa}: Se o professor informar valor negativo ou superior a 10,00, a restrição \texttt{chk\_grade\_range} do banco barra a gravação e a interface emite alerta.\\
& \textbf{[6e] Feedback Vazio}: O sistema exige preenchimento de ao menos uma justificativa pedagógica para garantir o caráter formativo da avaliação. \\
\textbf{Pós-condições} & Submissão avaliada com nota e parecer registrados no banco de dados. \\
\textbf{Regras de Negócio} & \textbf{RN13} (Escala de Nota de 0 a 10) e Transparência Pedagógica. \\
\textbf{Requisitos Associados} & \texttt{RF22}, \texttt{RNF03}. \\
\bottomrule
\end{tabularx}
\fonte{Elaborado pelo autor (2026).}
\end{table}

% ---------------------------------------------------------------------------
% UC20
% ---------------------------------------------------------------------------
\subsection{UC20: Comparar Gatilhos do Discente com Gabarito Docente}
\label{uc20_especificacao}

\begin{table}[!ht]
\centering
\caption{Especificação Detalhada do Caso de Uso UC20}
\label{tab_uc20}
\footnotesize
\begin{tabularx}{\textwidth}{p{3.5cm} X}
\toprule
\textbf{Propriedade} & \textbf{Especificação Técnica e Comportamental} \\
\midrule
\textbf{Identificador e Nome} & \texttt{UC20: Comparar Gatilhos do Discente com Gabarito Docente} \\
\textbf{Atores} & \texttt{PROFESSOR} (Ator Primário). \\
\textbf{Nível de Meta} & Nível Subfunção / Análise Comparativa (\textit{Fish Level}). \\
\textbf{Pré-condições} & Submissão discente aberta para correção no âmbito de \texttt{UC19}; atividade dotada de gabarito clínico de referência cadastrado pelo professor. \\
\textbf{Gatilho} & Invocado pelo fluxo principal de \texttt{UC19} durante o carregamento da submissão. \\
\textbf{Fluxo Principal} & 
1. O sistema recupera a lista de gatilhos enviada pelo aluno em \texttt{identified\_triggers}.\\
& 2. O sistema recupera a lista de gatilhos verdadeiros do gabarito docente.\\
& 3. O sistema calcula a matriz de confusão clínica:
\begin{itemize}
    \item \textbf{Verdadeiros Positivos (VP)}: Gatilhos corretamente identificados pelo aluno;
    \item \textbf{Falsos Positivos (FP)}: Gatilhos apontados pelo aluno que não existiam no caso;
    \item \textbf{Falsos Negativos (FN)}: Gatilhos reais existentes que foram omitidos pelo aluno.
\end{itemize}\\
& 4. O sistema calcula as métricas diagnósticas:
\begin{itemize}
    \item $\text{Sensibilidade} = \frac{VP}{VP + FN}$
    \item $\text{Concordância de Gravidade NCC MERP}$ (porcentagem de exatidão de categorias).
\end{itemize}\\
& 5. O sistema renderiza a interface comparativa lado a lado com destaques cromáticos: verde esmeralda (\texttt{\#059669}) para acertos, vermelho carmim (\texttt{\#DC2626}) para falsos positivos e cinza com alerta para omissões. \\
\textbf{Pós-condições} & Métricas comparativas exibidas ao professor para subsidiar a nota final. \\
\textbf{Regras de Negócio} & \textbf{RN14} (Algoritmo de Cálculo de Sensibilidade e Concordância). \\
\textbf{Requisitos Associados} & \texttt{RF23}, \texttt{RNF03}. \\
\bottomrule
\end{tabularx}
\fonte{Elaborado pelo autor (2026).}
\end{table}

% ---------------------------------------------------------------------------
% UC21
% ---------------------------------------------------------------------------
\subsection{UC21: Consolidar e Visualizar Indicadores Epidemiológicos IHI}
\label{uc21_especificacao}

\begin{table}[!ht]
\centering
\caption{Especificação Detalhada do Caso de Uso UC21}
\label{tab_uc21}
\footnotesize
\begin{tabularx}{\textwidth}{p{3.5cm} X}
\toprule
\textbf{Propriedade} & \textbf{Especificação Técnica e Comportamental} \\
\midrule
\textbf{Identificador e Nome} & \texttt{UC21: Consolidar e Visualizar Indicadores Epidemiológicos IHI} \\
\textbf{Atores} & \texttt{PROFESSOR}, \texttt{ADMIN} (Atores Primários). \\
\textbf{Nível de Meta} & Nível Usuário (\textit{Sea Level}). \\
\textbf{Pré-condições} & Existência de auditorias submetidas e consolidadas no banco de dados. \\
\textbf{Gatilho} & O docente ou administrador clica em "Painel Epidemiológico GTT". \\
\textbf{Fluxo Principal} & 
1. O usuário seleciona o escopo de análise: por turma específica, por período letivo ou global institucional.\\
& 2. O sistema executa agregação SQL/JSONB e calcula os indicadores canônicos do IHI \cite{ihi2009}:
\begin{itemize}
    \item Taxa de Eventos Adversos por 1.000 Pacientes-Dia;
    \item Taxa de Eventos Adversos por 100 Admissões Hospitalares;
    \item Porcentagem de Internações Hospitalares com Dano Real (\% com Danos E a I);
    \item Distribuição percentual de severidade de dano na escala NCC MERP (A a I);
    \item Ranking dos 10 gatilhos mais frequentemente disparados no período.
\end{itemize}\\
& 3. O sistema renderiza gráficos de rosca, barras empilhadas e cartões métricos dinâmicos com a paleta visual clínica suave da plataforma. \\
\textbf{Pós-condições} & Indicadores epidemiológicos apresentados em tela para análise pedagógica e institucional. \\
\textbf{Regras de Negócio} & \textbf{RN15} (Fórmulas Epidemiológicas do Institute for Healthcare Improvement). \\
\textbf{Requisitos Associados} & \texttt{RF24}, \texttt{RNF03}. \\
\bottomrule
\end{tabularx}
\fonte{Elaborado pelo autor (2026).}
\end{table}

% ---------------------------------------------------------------------------
% UC22
% ---------------------------------------------------------------------------
\subsection{UC22: Exportar Relatórios de Desempenho e Indicadores Clínicos}
\label{uc22_especificacao}

\begin{table}[!ht]
\centering
\caption{Especificação Detalhada do Caso de Uso UC22}
\label{tab_uc22}
\footnotesize
\begin{tabularx}{\textwidth}{p{3.5cm} X}
\toprule
\textbf{Propriedade} & \textbf{Especificação Técnica e Comportamental} \\
\midrule
\textbf{Identificador e Nome} & \texttt{UC22: Exportar Relatórios de Desempenho e Indicadores Clínicos} \\
\textbf{Atores} & \texttt{PROFESSOR}, \texttt{ADMIN} (Atores Primários). \\
\textbf{Nível de Meta} & Nível Usuário (\textit{Sea Level}). \\
\textbf{Pré-condições} & Visualização do painel de turmas ou painel epidemiológico. \\
\textbf{Gatilho} & O professor clica em "Exportar Relatório". \\
\textbf{Fluxo Principal} & 
1. O usuário escolhe o tipo de relatório desejado: Desempenho por Turma (notas, acertos e tempos de auditoria) ou Relatório Epidemiológico GTT.\\
& 2. O usuário seleciona o formato de saída: \texttt{CSV} para análise em pacotes estatísticos (R, SPSS) ou documento diagramado formatado para impressão.\\
& 3. O sistema compila o conjunto de dados aplicando os filtros vigentes.\\
& 4. O sistema gera o fluxo de bytes correspondente e dispara o download automático no navegador do cliente.\\
& 5. O sistema exibe notificação de exportação concluída. \\
\textbf{Pós-condições} & Arquivo de dados baixado na máquina do usuário sem alteração do banco de dados. \\
\textbf{Regras de Negócio} & Anonimização dos registros para garantia de conformidade com a LGPD. \\
\textbf{Requisitos Associados} & \texttt{RF25}, \texttt{RNF02}. \\
\bottomrule
\end{tabularx}
\fonte{Elaborado pelo autor (2026).}
\end{table}
